% arc_captions.tex -- GENERATED by figs_arc.py (do not edit by hand; rerun the script).
% Caption macros for figs/arc_{anatomy,small,skeleton,large,split,twocoset,gk2d,whylat}.pdf; parameters are those drawn.
\newcommand{\arcAnatomyCaption}{%
The mechanism of the nets of this paper, on a toy two-dimensional example ($q=9$; schematic, not a net from the text). (a) A $11\times11$ window of cells (grey grid), one point per cell. The geometric digit $r=Mu\bmod q$ is linear in the last cell digits $u$, with $M=-I+R_{90^\circ}$ (the $xy$-block of the geometric layer of the three-dimensional nets, Proposition~\ref{prop:cell}); it puts the point of cell $n$ at the lattice image $n+r/q$ (open circles). On each patch this is a piece of $(I+\frac Mq)\mathbb{Z}^2$, which is $\mathbb{Z}^2$ rotated by $\arctan\frac1{q-1}$ and contracted, with shortest vectors $(I+\frac Mq)e_x,(I+\frac Mq)e_y$ (arrows). Where a coordinate of $r$ wraps modulo $q$ the patch ends at a slip (dashed), which keeps one point per cell. The deeper digits add a jitter $f/q$, $f\in[0,1)^2$ (grey segments to the actual points, dots). Since $\det(I+\frac Mq)=(1-\frac1q)^2+\frac1{q^2}<1$, the lattice image is about $2/q$ too dense: about every $q$ cells one lattice site is skipped ($\times$) and the point jumps by about one cell at a slip; here 27 skipped sites against the density surplus $(\det^{-1}-1)\,11^2=29.7$, up to boundary effects. (b) The dotted box of (a): the jitter stays in a sub-cell of side $1/q$, so the separation is at least the lattice minimum $\lambda_1=0.895$ minus the jitter cost. (c) One axis, $r=-u\bmod q$ (the diagonal entry $-1$ of $M$): the points of cells $1,\dots,q$ sit on the sites $1+i(1-\frac1q)$, $i\in\Z$; cell $q$ receives a second site, which is skipped, and the next point steps $2-\frac1q$ across the slip. Theorem~\ref{thm:sep} makes this exact ($1-1/q$ from the own-axis step); the other regimes (II and III, Section~\ref{sec:prelim}) change only the lattice and the cost.}
\newcommand{\arcSmallCaption}{%
Regime I, small-entry layers ($m=sk$). (a) $\mathcal P_{2,11}$ ($q=11$, $k=2$, $N=11^6$; Theorem~\ref{thm:main}): the 196 points with $z$-cell $n_z=60$ and $(x,y)$-cells in a $14\times14$ window, projected to the $xy$-plane in cell units. Open circles: the lattice image $n+r/q$, with $r=Mu+Jv\bmod q$ the geometric digit of Proposition~\ref{prop:cell}; dots: the points; grey segments: the jitter ($<1/q$ per coordinate). Arrows: the $(x,y)$ parts of the shortest vectors $(I+\frac Mq)e_x$, $(I+\frac Mq)e_y$ of $\Lambda_q=(I+\frac Mq)\mathbb{Z}^3$, $\lambda_1(\Lambda_q)=0.918$ (Remark~\ref{rem:lattice3}). Dashed: slip edges, where a geometric digit wraps; dotted: cell boundaries where the last cell digit wraps (the term $Jv$ enters); $\times$: skipped sites. As in Figure~\ref{fig:arc}, $\Lambda_q$ is too dense, $\det(I+\frac Mq)=0.773\approx1-\frac3q$: 35 skipped sites against the density surplus 38.8, up to boundary effects. (b) Dimension four: a $(0,8,4)$-net over $\mathbb F_{17}$ ($k=2$, $N=17^8$) whose row~$0$ is the deficit-two design over $\Z[1/2]$ (places $\infty,0,1,2$; Section~\ref{sec:dim4ex}), $xy$-slice with $n_z=144$, $n_w=96$ fixed; arrows: two of the eight shortest vectors $\pm(I+\frac{F_0}q)e_c$ of $(I+\frac{F_0}q)\mathbb{Z}^4$ ($\lambda_1=0.888$). Despite the larger base, the entries of size two of $F_0$ (Remark~\ref{rem:pm1}) make the patches smaller and the slips denser than in (a): 68 against 58 slip edges in the same window, and 46 $\times$ against a planar density surplus of 53.5, up to boundary effects. The separation bound of Section~\ref{sec:dim4ex} uses only row~$0$.}
\newcommand{\arcSkeletonCaption}{%
Regime II, lattice skeletons ($m=sk+r$), in fine-cell units. (a) $\mathcal B^2_{1,43,3}$ ($p=43$, $\rho=3$, $m=5$; Remark~\ref{rem:B2}): the slab of fine cells $P_z=1000$ over a $30\times30$ window. Shaded: the occupied fine cells, which are exactly the cells of the index-$p$ lattice $\Lambda_{\mathcal B^2}=\{a:a_z\equiv\rho a_x+\rho^2a_y\pmod p\}$: no faults, no skipped sites. Open circles: their corners (the lattice image); dots: the points, displaced by up to one fine cell of jitter; arrows: a reduced basis $(-3,1)$, $(4,13)$ of the lattice section. (b) The three slabs $P_z=1000,1001,1002$ (shades): the squares of side $\rho$ about the occupied cells have disjoint interiors, so two points of these slabs differ by at least $\rho$ fine cells in $x$ or $y$ (Lemma~\ref{lem:skellat}(L)), and $\delta N^{1/3}\ge(\rho-1)p^{-1/3}$ after one fine cell of jitter. (c) $\mathcal B^1_{1,37,3}$ ($m=4$; Theorem~\ref{thm:allm}): $xz$-view of the points with $y$-cell index $5$ over the four coarse blocks around $(37,37)$ (black lines). Open circles: the coset of the Korobov lattice $\Lambda_{\mathcal B^1}$ (generating vector $(1,\rho,\rho^2)$ modulo $p$), continued from the lower-left block (projected basis $(-4,1)$, $(1,9)$); dashed arrows: the faults. In each coarse block the occupied cells are exactly a Korobov coset, shifted by $\pm1$ fine cell from block to block (dark where shifted): the faults are coset shifts, like stacking faults, not skips. The separation is the shortest vector minus faults and jitter (Lemma~\ref{lem:skelred}).}
\newcommand{\arcLargeCaption}{%
Regime III, large-entry layers ($m=s$). (a) The column net $\mathcal C_{13}$ (Theorem~\ref{thm:col}, $a=(4,4,5)$, $w=(-1,-1,2)$): the $3^3$ cells around $n=(8,8,3)$. All three coordinates read the same two offset digits, so each point lies exactly on the body diagonal of its cell (dotted), at the fraction $\varphi(n)/p^2$ of it, which changes by $\frac13+\frac{w\cdot E}{3p}$ from cell to cell, up to the jitter $\Delta j/p^2$. Open circles: the reflected cube $H\mathbb{Z}^3$, $H=I-\frac23\mathbf 1\mathbf 1^{\!\top}$, anchored at the central point; accent: its six minimal vectors $\pm He_c$; grey segments: drift and jitter, both along $(1,1,1)$. (b) The same net, the $5^3$ cells around $n$, viewed along $(1,1,1)$: the lattice points within one plane spacing $1/\sqrt3$ of the centre along $(1,1,1)$ (fill: the plane $x+y+z=\mathrm{const}$). The three neighbours one plane up, $-He_c$, project opposite to the cell axes $e_c$ (grey), whereas in $\mathbb{Z}^3$ they are the $e_c$ themselves: $H\mathbb{Z}^3$ is $\mathbb{Z}^3$ turned half a turn about $(1,1,1)$. Drift and jitter run along $(1,1,1)$ and are invisible here: every drawn point projects exactly onto its site, and differs from its site of the anchored lattice by less than $0.15$ of a cell in each coordinate. (c) Dimension five: the net $\mathcal H^{5*}_{1,p}$ of Proposition~\ref{prop:HL5} with $k=1$ ($m=s=5$), $p=10007$, a $(0,5,5)$-net over $\mathbb F_p$: the $xy$-slice of an $8\times8$ window of cells, the other three cells fixed. Its layer $F_0$ lies in $\frac12+\mathbb{Z}$, so by the parity of $\Sigma n$ the point sits at $n+c_0$ or, in the shaded cells, at $n+c_0+\frac12\mathbf 1$, for one fixed $c_0\in[0,1)^5$ (here $c_0\approx0.10$ in $x$ and $y$), up to drift and jitter below $0.04$ of a cell: locally the packing $D_5^+$, whose minimum $\sqrt5/2>1$ (a step to a cell of the other parity) exceeds the cubic value~$1$; a diagonal step within one parity has length $\sqrt2$ (Proposition~\ref{prop:Hk} with $k=1$). In all three panels the limit lattice has covolume~$1$ and one site in each drawn cell: no site is skipped, and the net departs from the lattice only by the $O(1/p)$ drift and jitter.}
\newcommand{\arcSplitCaption}{%
Split frames (Definition~\ref{def:split}). (a) Top: a frame reads the cell and the offset of coordinate $c$ at one place $a_c$, so its geometric layer is a symmetric Cauchy matrix, and a Pl\"ucker relation among four places rules out a half-integer layer (Proposition~\ref{prop:P}). Bottom: the split frame of $H_5^*$ (Proposition~\ref{prop:HL5}), $u=(-2,0,2,4,8)$ even, $v=(5,7,-1,3,1)$ odd: coordinate $c$ reads its cell at $v_c$ (filled) and its offset at $u_c$ (open); arcs: two entries $\vartheta_{cd}=u_c-v_d$. The box minors are confluent Vandermonde determinants in the ten places (Lemma~\ref{lem:Tsplit}). (b) The geometric layer is the non-symmetric Cauchy matrix $\lambda^\prime_c\mu_d/(u_c-v_d)$ (Proposition~\ref{prop:Gsplit}); every $\vartheta_{cd}$ is odd, so the layer can be put in $\frac12+\mathbb{Z}$ (Lemma~\ref{lem:Hs}; for the explicit $H_s$, $u_c=2c$, $v_d=2d+1$). Shown: $\vartheta$ and $2F_0$ of $H_5^*$. (c) The net $\mathcal H^{5*}_{2,p}$, $p=10007$: the $xy$-slice of a $10\times10$ window of cells that crosses a carry in $x$ and in $y$ (dotted), the other three cells fixed. Every partial sum $M^{(j)}$ of the design lies in $\frac12+\mathbb{Z}$, so the corner/centre parity pattern of $D_5^+$ (minimum $\sqrt5/2>1$) continues across the carries (Proposition~\ref{prop:Hk}); drift and jitter stay below $0.05$ of a cell.}
\newcommand{\arcTwoCosetCaption}{%
A two-coset layer at $m=3$: the net $\mathcal T_{1009}$ of Proposition~\ref{prop:twocoset} ($h=504$, $r=126$), in cell units. (a) The $(y,z)$-plane of the slab of cells $n_x=661$, $7\times7$ cells around $(n_y,n_z)=(309,171)$. The coordinates $y$ and $z$ read the same geometric digit $rn_x-n_y+n_z$, so the points of a slab lie, up to drift and jitter, on a translate of $\mathbb{Z}^2$ (open circles: the pattern $\Gamma$ anchored at the central point); the geometric digit $h(n_y+n_z)$ of $x$, with $h\equiv-\frac12$, raises the points of every other cell (shaded) by $\frac12$ in $x$. Grey dots: the next slab $n_x+1$, shifted by $(\frac18,\frac18)$. Accent: from a raised point, the two kinds of shortest vectors to the next slab, $(1,\frac18,\frac18)$ and $(\frac12,-\frac78,\frac18)$, both of length $\sqrt{66}/8$. (b) The $(x,y)$-plane of the cells with $n_z=171$: along $y$ the points alternate between the two heights (the pucker; in-slab neighbours are $\sqrt5/2$ apart), and from column to column the $y$-offset grows by $\frac18$ (the shear) and wraps after eight columns. Accent: the same two shortest vectors (their $z$-components $\frac18$ are not visible) and an in-slab neighbour. In (a) and (b) every point is within $0.01$ of its site of $\Gamma$ in each coordinate, and $\Gamma$ has one site in each cell; it is the union of two cosets of a lattice of covolume $2$ and is not a lattice. (c) $\delta N^{1/3}$ against $p$: exact values for $\mathcal T_p$ at 8 primes $101\le p\le1009$ (dots; computed, Appendix~\ref{app:verify}), the lower bound of Proposition~\ref{prop:twocoset}(2) (for $p\ge41$), the pair value $d_*/p^2$ of (3), and for comparison the exact values of the column nets $\mathcal C_p$ (Theorem~\ref{thm:col}; solid for $p\equiv2$, dotted for $p\equiv1\pmod3$), which stay below $1$ (Corollary~\ref{cor:colceil}); dashed: $1$ and $\sqrt{66}/8$.}
\newcommand{\arcGKCaption}{%
The mechanism in its simplest instance: the permutation $(0,m,2)$-net of Gr\"unschlo{\ss} and Keller \cite[Sec.~3]{GK09} for $m=2r=8$, $Q=2^r=16$, $N=Q^2=256$, in cell units (Remark~\ref{rem:gk2d}). (a) The $(r,r)$ cells (grey grid), one point each. Open circles: the lattice image $ue+jv+(\frac14,0)$, $e=(1,0)$, $v=(\frac14,1)$, on the lines (dashed) of $\Lambda=\langle e,v\rangle$ in direction $v$; the tiling lattice of the row-shifted torus is $Q\Lambda$, so the lines close up and there are no slips. Dots: the points $(X,Y)=ue+(j+\varrho(u)/Q)v+(d,0)$, $|d|<1$ grid unit: line $u$ is slid along itself by its slide $\varrho(u)/Q$ of one step, a van der Corput value (grey segments, the jitter). Accent: the basis $e$ (shortest, length $1$) and $v$. (b) Zoom on $6\times6$ cells with the cells of depths $(r+2,r)$ (light grid); the occupied ones (shaded) are exactly the coset $\{(a,b):a\equiv b+1\pmod 4\}$ of the index-$4$ lattice $\Lambda$. Since the slides move the points arbitrarily along the lines, only the line spacing $4/\sqrt{17}$ (double arrow) survives as a separation bound: $\delta\sqrt N\ge\frac4{\sqrt{17}}(1-\frac1Q)=0.909$ for every choice of slides; the van der Corput slides give exactly $\delta^2N=241/256$ (in the tiling metric; \cite{Rud26}).}
\newcommand{\arcWhyLatCaption}{%
What linearity forces (Section~\ref{sec:whylat}). (a) Two bijections of $\Z/W$, $W=24$, as sub-cell offsets along a cyclic line of cells: the affine bijection $i\mapsto i$ (grey), the kind that a linear offset gives, jumps once, by $W-1$, and this jump sits at a slip plane (Corollary~\ref{cor:linoff}); every affine bijection of $\Z/W$ has a step of size at least $W/2$, since its steps are congruent modulo $W$ and sum to zero. A tent-shaped bijection (accent), available only to nets that are not digital, moves by at most $2$ between neighbouring cells (Remark~\ref{rem:open}(11)). (b) The wrap vectors $\omega$ of Lemma~\ref{lem:phase}(i) that occur for the step $E=\e_1$ over all $p^3$ cells, $p=31$ (filled), for the equal-row layer $\mathbf 1a^{\!\top}$, $a=(10,10,11)$, of the column net $\mathcal C_{31}$ (Definition~\ref{def:col}) and for the invertible layer obtained by adding $\left(\begin{smallmatrix}0&1&0\\0&0&2\\1&0&0\end{smallmatrix}\right)$; framed, the pattern of least length $|E+b-\omega|$, printed on the right (computed: $1.010$ and $0.577$). The shared wrap allows only $000$ and $111$; the perturbed layer realizes all eight, and with the glue $b\approx\frac13\mathbf 1$ (Lemma~\ref{lem:phase}) the fully folded length drops to $1/\sqrt3$. (c) A coordinate that does not move contributes $\min(b_c,1-b_c)$ under free wraps, at most $\frac12$, attained only at $b_c=\frac12$ (Remark~\ref{rem:glue}).}
